\SourceLeaf{169}the sum of all the approaches, however many times they were continued,
would always remain less than $\frac{1}{1-\psi} \times ckf$, and
this bound could still be less than the right angle
$dk\varphi$. You asked for my candid judgment. I
have given it, and I assure you once again that
I will be truly glad if you overcome all the difficulties.

.\, .\, .\, .\, .\, .\, .\, .\, .\, .\, .\, .\, .\, .\, .\, .\, .

\medskip
\begin{center}
\rule{0.3\textwidth}{0.4pt}
\end{center}
\medskip

\begin{center}
\textsc{Remarks.}
\end{center}

\textsc{Gauss} and \textsc{Wolfgang Bolyai} of Bolya in
Transylvania studied together in Göttingen from autumn 1796 to autumn 1798
and were friends. Before
\textsc{Bolyai}'s return home they met once more, on 25 May 1799 in
Clausthal. The letter of 16 December 1799, an extract of
which is printed here, is the first that \textsc{Gauss}
addressed to his friend after they parted.

\textsc{Bolyai} sent his »Göttingische Theorie der Parallelen« to \textsc{Gauss} only in 1804.
The letter of 25 November 1804 contains Gauss's critique.

The following entries supplement the foregoing statements about \textsc{Gauss}'s
investigations into the first principles of
geometry.

\begin{quote}
\emph{»Plani possibilitatem demonstravi. [1797] Jul. 28. Gotting.«}
\end{quote}
and
\begin{quote}
\emph{»In principiis Geometriae egregios progressus fecimus.
Br[unovici 1799] Sept.«}
\end{quote}
They are found in the \textls{diary} already mentioned
on p. 20 of this volume.

\begin{flushright}
\textsc{Stäckel.}
\end{flushright}

\medskip
\begin{center}
\rule{0.3\textwidth}{0.4pt}
\end{center}
